\chapter{Testing and Quality Assurance}
\label{ch:testing}

This chapter documents the test strategy that supports the requirements in \Cref{ch:requirements}. APEX is tested at three levels: backend unit and integration tests with the standard Go testing toolchain, frontend component tests with Vitest, and manual end-to-end exercises against the running app. Static checks (\texttt{go vet}, ESLint, TypeScript) catch the kind of mistakes that tests should not have to.

\section{What is tested, and where}

\Cref{tab:test-pyramid} summarises the surface covered at each level. The split is deliberate: anything that depends on PostgreSQL or Judge0 is best tested as an integration test, not a unit test, because the grading and finalisation logic is dominated by the database queries and HTTP round trips.

\begin{table}[H]
  \centering
  \caption{Test surface by level.}
  \label{tab:test-pyramid}
  \small
  \begin{tabularx}{\linewidth}{l X X}
    \toprule
    Level & Coverage focus & Tooling \\
    \midrule
    Unit (Go)            & pure logic in \texttt{judge0} (output normalisation, MCQ set comparison), JWT claim extraction, time parsing, route registration  & \texttt{go test} \\
    Integration (Go)     & handler+database flows: signup-login, attempt lifecycle, finalisation, manual grade overrides, IDOR rejection, password reset round-trip & \texttt{go test} with a live Postgres test DB  \\
    Unit (frontend)      & utility functions, component rendering with stubbed data (Question card, AnswerEditor, Timer)                                          & Vitest, \texttt{frontend/src/test/setup.ts}  \\
    End-to-end (manual)  & golden-path: build $\to$ assign $\to$ take $\to$ grade $\to$ review for one coding, one MCQ, one written exam                  & \texttt{./start.sh} with seeded data           \\
    Static               & type checks, linting, vet                                                                                                              & \texttt{tsc}, \texttt{eslint}, \texttt{go vet} \\
    \bottomrule
  \end{tabularx}
\end{table}

\section{Backend tests}

The backend tests live alongside their packages in \texttt{*\_test.go} files. The convention follows Go's standard testing toolchain~\cite{go}: each test exercises one exported function or one HTTP handler, sets up its own fixtures, and tears them down explicitly.

\paragraph{Unit tests.} The hardest-to-write logic in the system is also the easiest to unit-test. \texttt{judge0.normalizeOutput} (\texttt{internal/judge0/grader.go:366--374}) has a small table-driven test that asserts: trailing whitespace is stripped per line, \texttt{\textbackslash r\textbackslash n} collapses to \texttt{\textbackslash n}, and edge whitespace is trimmed. The MCQ set-equality check inside \texttt{judge0.GradeMCQ} is unit-tested with a fake DB shim. \texttt{exam.parseTime} (\texttt{internal/exam/handler.go:737--751}) is unit-tested across all five accepted layouts.

\paragraph{Integration tests.} The integration suite uses a dedicated Postgres test database (\texttt{apex\_test}) brought up by Docker Compose. Each test starts a transaction, runs its scenario, and rolls back so tests are independent. The flows under test are:

\begin{itemize}
  \item Signup, login, login with bad password, login on a Google-only account (empty hash).
  \item Forgot-password $\to$ reset-password $\to$ login-with-new-password, including the stale-token rejection path.
  \item Create class, join with invite code, leave, remove member.
  \item Build a draft exam, fail to publish without \texttt{start\_time}, set start time, publish, assign.
  \item Start an attempt twice (idempotent), submit, second submit returns 409.
  \item Auto-grade a coding submission against a stubbed Judge0 (an HTTP test server returning canned responses).
  \item Manually grade a written submission, override, and verify the attempt's aggregate score updates.
  \item Reject IDOR on \texttt{POST /submissions/run} (the F1 fix from \Cref{ch:security}).
\end{itemize}

\paragraph{Stubbing Judge0.} The auto-grader's only HTTP client lives in \texttt{judge0.RunCode} and uses \texttt{judge0.BaseURL} as the base. Integration tests spin up an \texttt{httptest.Server} that serves canned \texttt{/submissions} responses, then re-point \texttt{BaseURL} at the test server's URL via \texttt{judge0.Init}. This keeps the integration tests deterministic without losing coverage of the surrounding HTTP marshalling.

\section{Frontend tests}

The frontend test setup lives in \texttt{frontend/src/test/setup.ts} and is wired into Vitest via \texttt{vitest.config} (inherited from \texttt{vite.config.ts}). The seed test, \texttt{frontend/src/test/example.test.ts}, asserts that the test runner is wired correctly.

\paragraph{Component tests.} Component tests use \texttt{@testing-library/react} to render a component with stubbed props and assert visible text and interactions. Three components have current coverage:

\begin{itemize}
  \item The Markdown renderer in \texttt{ExamTaking.tsx} is asserted not to render raw HTML when its source contains \texttt{<script>} or \texttt{<img onerror=...>}.
  \item The autosave timer hook stubs \texttt{api.autosaveExamAttempt} and asserts the call cadence.
  \item The MCQ option list asserts that selecting an option toggles its underlying ID in the answers array.
\end{itemize}

\paragraph{What is intentionally not unit-tested.} Page-level components that depend on routing, auth context, and the API client are exercised by manual end-to-end runs rather than by component tests. The reasoning is the standard one: a component test that mocks half the world tells you mostly that the mocks are consistent with each other.

\section{Manual end-to-end exercises}

The end-to-end test plan is run at the close of every implementation phase and against every release candidate. The seed data set creates two users and one class:

\begin{itemize}
  \item \textbf{Dr.\ Demo} (\texttt{teacher@demo.test}, role \texttt{teacher}).
  \item \textbf{Sam Student} (\texttt{student@demo.test}, role \texttt{student}).
  \item One class (\texttt{Demo CS101}, invite code \texttt{ABCD1234}) with Sam enrolled.
  \item One coding exam (\texttt{Sum Two Numbers}) with three test cases (one sample, two hidden).
  \item One MCQ exam (\texttt{Time Complexity Quiz}) with four questions.
  \item One written exam (\texttt{Discuss Recursion}) with one question and a rubric.
\end{itemize}

The golden path exercises:

\begin{enumerate}
  \item Dr.\ Demo signs up, creates the class, builds and assigns each exam.
  \item Sam signs up, joins the class with the invite code, sees the assigned exams.
  \item Sam takes the coding exam: edits in Monaco, runs samples, submits. Dr.\ Demo's results screen shows the auto-graded score and per-test detail.
  \item Sam takes the MCQ exam, submits. Auto-grader marks it accepted/wrong-answer in milliseconds.
  \item Sam takes the written exam, submits. Dr.\ Demo's pending-grading queue shows one item; she grades it with feedback. Sam receives the "Exam Graded" notification.
  \item Dr.\ Demo toggles \texttt{grades\_announced}; Sam's review screen reveals the aggregate score and feedback.
\end{enumerate}

\paragraph{Edge-case exercises.} The plan also covers: take-after-end-time (rejected with 403), publish-without-schedule (rejected with 400), draft autosave $\to$ refresh page $\to$ resume on the same device, draft autosave $\to$ resume on a fresh browser (server-side draft is preferred), submit twice (second is 409), grade an attempt then override and re-grade (aggregate updates), delete a class with active attempts (cascades cleanly).

\section{Static checks}

\paragraph{Go.} \texttt{go vet ./...} runs as part of the CI step and catches printf-mismatch, shadowed return values, lock-by-value, and the other patterns vet knows about. The codebase compiles without warnings.

\paragraph{TypeScript.} \texttt{npm run lint} runs ESLint with the React and TypeScript plugins. The build (\texttt{npm run build}) refuses to emit if \texttt{tsc} reports an error, so type errors are surfaced before they reach a deploy.

\paragraph{Repository conventions.} Two markdown documents in the repo root encode rules that the type system cannot: \texttt{ORM\_RULES.md} for migration patterns, and \texttt{BUG\_FIX\_PLAN.md} for the running list of phases and their issue traces. New work is expected to update both.

\section{Coverage and gaps}

Coverage is measured by \texttt{go test -cover} and \texttt{vitest --coverage}. The project does not enforce a minimum percentage --- the experience is that headline coverage numbers are easy to manipulate and not the metric a maintainer cares about --- but a few subsystems are deliberately well-covered: \texttt{judge0} (output normalisation and the three graders), \texttt{exam.SubmitAttempt}, \texttt{auth.Signup}/\texttt{Login}, and the password-reset round trip.

The known gaps are:

\begin{itemize}
  \item The reminder scheduler is exercised manually by setting an exam's \texttt{start\_time} 50 minutes in the future and watching the email log; an integration test that wraps a fake clock would close the gap.
  \item The Google OAuth path is exercised manually because the Google SDK refuses to validate test tokens. A handler-level test that injects a fake \texttt{idtoken.Validate} via a function variable would let it run on CI.
  \item The end-to-end test plan is manual; an automated Playwright suite is on the roadmap (\Cref{ch:conclusion}).
\end{itemize}

\section{Summary}

The test surface is shaped by what the code can lie about. The pieces that compute (output normalisation, MCQ set equality, finalisation arithmetic, JWT parsing) are unit-tested. The pieces that travel through the database (attempt lifecycle, ownership checks, password-reset round-trip) are integration-tested. The pieces that are dominated by user judgment (the take screen, the grading queue, the results review) are manually tested against a small seeded data set. The result is a suite that runs in seconds on every push and a manual checklist that runs in about ten minutes per release candidate.
